%----------------------------------------------------------------------
\section{Discussion and Conclusion}\label{sec:Conclusion and Discussion}
%----------------------------------------------------------------------

% Formal verification of quantum error correction is a necessary component of
% trustworthy quantum computation. We have presented a Lean~4 framework that
% combines interactive theorem proving with SAT-based computation to verify
% nontrivial code properties, including distance claims for practical CSS and
% Bivariate Bicycle code instances. The framework includes a reusable formal
% library for stabilizer-code theory, a verified translation from code-distance
% problems to Boolean satisfiability, and demonstrated verification of several
% concrete codes.
%
% The result is a stronger trust story than either testing-only or solver-only
% workflows: every step from quantum-theoretic definitions through matrix-level
% computations to final distance claims is checked by Lean's proof kernel. As
% quantum hardware matures and fault-tolerance becomes critical, such formally
% verified foundations will play an increasingly important role in ensuring the
% reliability of quantum computations.
%

We presented \textsc{Lean-QEC}, a Lean 4 framework that combines interactive
theorem proving with SAT-based computation to deliver end-to-end,
machine-checked distance certificates for practical CSS and Bivariate Bicycle
codes. The library contributes a reusable stabilizer-code stack, a verified
reduction from code distance to Boolean satisfiability, two scalability
techniques (\lstinline{BitVec} flattening and location-indexed error encoding),
and case studies covering the Steane, Golay, and a family of Bivariate Bicycle codes up to $[[90, 8, 10]]$ inside
the kernel and the $[[144,12,12]]$ instance outside it.

The result narrows the gap that has separated the testing-only and solver-only
options for QEC verification. Every step from the quantum-theoretic definition
of distance, through the binary symplectic representation, through the
\lstinline{BitVec} encoding, to the final SAT certificate is connected by a
Lean-checked theorem. Either the kernel accepts the artifact or it does not; no
informal mathematical glue sits in between.

\subsection{Methodological Takeaways}\label{sec:method}

% The key contribution of this work is methodological: an end-to-end,
% proof-aware verification path that encompasses both abstract QEC theory and
% computationally intensive property checks. This narrows the gap between
% mathematically sound reasoning and practical verification for real code
% families.
%
%
%
% The formalization effort revealed several design insights:
%
% \begin{itemize}[leftmargin=1.5em]
%   \item \textbf{Bridging Representations.} The binary symplectic
% representation serves as a critical bridge between the quantum-theoretic
% notion of stabilizer codes and a cleaner, more algebraically manipulable
% layer. The binary symplectic representation is further bridged by a boolean
% representation which allows compatibility with SAT solving. The formal
% verification of these bridges allows for greater end-to-end confidence in
% results.
%   \item \textbf{Computer Algebra.} The SAT pipeline requires explicit kernel
% bases as input. Computing these kernels is itself a linear-algebra task that
% must be verified or trusted. Our current approach pre-computes kernels
% externally and includes them as definitions, with their kernel property stated
% as a lemma to be verified. Lean largely lacks preexisting capability for
% verified algebraic computations, which emerges as an interesting direction for
% future research if Lean is to become more widely useful.
% \end{itemize}
%

The build process surfaced two design lessons that we expect to transfer to
future quantum formalizations.

\paragraph{Representation bridges are first-class objects.} The binary
symplectic representation is the obvious algebraic bridge between the
quantum-theoretic stabilizer code and the matrix-level SAT query. The
non-obvious step is that this bridge must itself be a formalized object, not an
informal correspondence, if the resulting certificate is to be end-to-end. We
further bridged the BSR to a \lstinline{BitVec} representation tuned for the
solver. Each bridge is one Lean theorem and, once written, is invisible at the
call site; without it, every downstream lemma would carry an unverified
soundness gap.

\paragraph{Verified linear algebra is the missing piece.} The pipeline still
admits one externally computed input: a kernel basis for each parity-check
matrix. We discharge its correctness via a rank/orthogonality argument, but the
kernel itself is supplied as a definition rather than computed. A verified
Gaussian-elimination implementation in Lean would let the framework take
parity-check matrices alone as input. More broadly, the absence of efficient
verified algebraic computation is the consistent friction point in this
development; closing it would benefit many Lean formalizations beyond QEC.

\subsection{Related Work}\label{sec:related}

% Formal verification of quantum computing has seen growing interest. Efforts in
% Coq and other proof assistants have addressed quantum circuit verification and
% quantum information theory. The SQIR framework~\cite{hietala2021} provides
% verified quantum circuit optimization in Coq. QHLProver~\cite{liu2019}
% formalizes quantum Hoare logic. On the Lean side, projects such as
% Lean-QuantumInfo have begun formalizing quantum information primitives
% \cite{meiburg2025formalization}. Other work has verified properties of
% Error-Corrected programs, but outside of Interactive Theorem Provers
% \cite{huang2025efficient, wu2021qecv}. Our work is distinguished by its use of
% an Interactive Theorem Prover, focus on the error-correction layer
% specifically, and by the integration of SAT solving for computational
% verification tasks that resist concise hand proofs. Additionally, our
% end-to-end approach gives us more confidence in our results, which we argue
% are more completely formalized than other similar efforts.
%
% \Ethan{TODO cite concurrent work.
% Even if they don't have paper yet, we can cite github.}
%

\paragraph{Quantum verification in proof assistants.} The SQIR
development~\cite{hietala2021} and CoqQ~\cite{coqq} formalize quantum circuits
and circuit optimizations in Rocq, and Peng et
al.~\cite{peng2022formallycertifiedendtoendimplementation} verify end-to-end
correctness of algorithms such as Shor's. QHLProver~\cite{liu2019} formalizes
quantum Hoare logic. On the Lean side,
Lean-QuantumInfo~\cite{meiburg2025formalization} formalizes finite-dimensional
quantum-information primitives but does not specialize to qubit systems or QEC;
we build on the same Mathlib foundations but specialize to the qubit space and
develop the stabilizer-code stack.

\paragraph{Verification of QEC.} The 9-qubit Shor-code verification
of~\cite{feng-9-qubits} closes the trust loop on a single small code but
enumerates Pauli errors explicitly and does not scale beyond it. Other
verified-QEC efforts~\cite{huang2025efficient, wu2021qecv} sit outside ITPs and
target program-level properties rather than code-distance certification. To our
knowledge, \textsc{Lean-QEC} is the first development that combines a
kernel-checked stabilizer stack with a solver-backed distance pipeline that
scales to industrial code sizes.

\paragraph{Solver-assisted ITP.} The strategy of reducing a hard subgoal to SAT,
dispatching it to an external solver, and reconstructing the certificate inside
the kernel is well established in classical verification. Our contribution to
this idiom is the verified reduction from a quantum-theoretic distance statement
to a SAT formula, the \lstinline{BitVec}-flattened encoding that makes the
formula tractable to instantiate, and the location-indexed variable scheme that
makes the formula tractable to solve.

Concurrent work by Jain~\cite{qeclean} also explores formalization of quantum
error-correcting codes in Lean but focuses on the toric code family.

\subsection{Limitations and Future Work}\label{sec:future}

\paragraph{Solver scaling.}
The SAT pipeline scales by orders of magnitude over hand-written ITP distance
proofs --- our 144-qubit instance handles roughly $2^{40}$ candidate errors ---
but it is not unbounded. Codes substantially larger than the BB family
($\mathrm{BB}_{360}$ and beyond) push the SAT query itself out of reach.
Decomposition strategies and specialized solvers for
linear codes are natural directions for extending the envelope.

\paragraph{Broader code families.}
The development targets CSS codes, with BB as the headline family. Extending to
general non-CSS stabilizer codes is mostly a matter of dropping the $X$/$Z$
split and reusing the full BSR pipeline; subsystem codes and codes from the
broader stabilizer formalism require more substantial additions. Note that
surface and color codes are themselves CSS and should fit the existing framework
directly.

\paragraph{Verified kernel computation.}
The kernel basis for each parity-check matrix is currently produced by an
external Python script and certified inside Lean. A verified row-reduction
routine would close this remaining preprocessor step. Doing so cleanly is the
most concrete piece of follow-on engineering implied by this work.

\paragraph{Integration with quantum-circuit verification.}
The QEC layer is one component of the verified quantum-computation stack.
Connecting it to verified circuit compilers such as SQIR would yield an
integrated, end-to-end pipeline from algorithm specification to fault-tolerant
physical implementation.
